\section{Конструкции, используемые в групповых выводах}\label{sec:constructions}

\subsection{Поле и базовая координатная пара}
Положим
\[
 K=\Q[a]/(a^2-a+1),\quad b=1-a,\quad ab=1,\quad b^2=-a.
\]
Чебышёвские полиномы в целочисленной нормировке определяются рекурсией
$D_0=2$, $D_1=t$, $D_{n+1}=tD_n-D_{n-1}$. В характеристике нуль
$D_n(t)=2T_n(t/2)$. В частности,
\[
 D_3=t^3-3t,\quad D_4=t^4-4t^2+2,\quad D_5=t^5-5t^3+5t,
\quad D_7=t^7-7t^5+14t^3-7t.
\]
Рассматриваем конкретные параметризации
\[
 \gamma_T=(D_3,D_4-2,t^5-10t),\qquad
 \gamma_H=(D_3,D_5,D_7).
\]
Во второй компоненте $\gamma_H$ коэффициент при $t$ равен $5$.
Параметризация с коэффициентом $4$ данным результатом не охватывается.

Аффинная карта $M=(x,y+2,z-5x)$ переводит $\gamma_T$ в
$\gamma_0=(D_3,D_4,D_5)$. В нормализованных координатах введём
\begin{align*}
 u&=axy+z,&v&=-y,&m&=bx^5,&s&=bx,\\
 \ell&=-1+ax^2-ax^4,&
 g&=-2x^5+(6-2a)x^3-2bx,\\
 P_x(U)&=sU^2+\ell U+g,&H&=P_x(u)+mv.
\end{align*}
В кольце $K[x,T]/(x^5)$ полином $P_x$ обратим. Используем представитель
\begin{align}
 F_x(T)={}&-T+x(bT^2+2a-2)+x^2(2aT^3-5aT)\notag\\
 &+x^3(-5T^4+15T^2-2a)\notag\\
 &+x^4\bigl(14bT^5+50(a-1)T^3+(29-24a)T\bigr).
 \label{eq:inverse-mod-five}
\end{align}
Прямое вычисление даёт обе конгруэнции
\[
 P_x(F_x(T))\equiv T\pmod{x^5},\qquad
 F_x(P_x(U))\equiv U\pmod{x^5}.
\]
Поэтому
\[
 W=\frac{u-F_x(H)}m,\qquad
 \chi_x(T)=\frac{T-P_x(F_x(T))}m
\]
являются полиномами. Для $A_0=(x,H,W)$ обратная карта имеет вид
\begin{align}
 U&=F_x(T)+mV,\notag\\
 V_0&=\chi_x(T)-(\ell+2sF_x(T))V-smV^2,\notag\\
 A_0^{-1}(x,T,V)&=(x,-V_0,U+axV_0).\label{eq:base-inverse}
\end{align}
Действительно, раскрытие квадрата даёт $P_x(U)+mV_0=T$ и
$(U-F_x(T))/m=V$. Для обратной композиции сначала восстанавливаются
$u,v$, затем $y,z$. Это полиномиальное доказательство обратимости;
постоянство якобиана само по себе здесь не используется как критерий.

Подстановки в кривую дают
\[
 H(D_3,D_4,D_5)=t,\qquad W(D_3,D_4,D_5)=8-14t^2.
\]
Следовательно, при $h_T=H\circ M$, $w_T=W\circ M$ получаем
\begin{equation}\label{eq:trefoil-rectifier}
 \Phi_T=(h_T,\ x-h_T^3+3h_T,\ w_T-8+14h_T^2),
 \qquad \Phi_T\gamma_T=(t,0,0).
\end{equation}
Обратное строится из \eqref{eq:base-inverse}, $M^{-1}$ и обратных
внешних сдвигов. Степени компонент данной карты равны $(5,15,24)$,
а якобиан равен $1$. Оптимальность степеней не заявляется.
Исходная реализация находится в \texttt{verify.py} репозитория;
новый проверяющий код отдельно воспроизводит компактные сертификаты.

\subsection{Перенос с контролем обеих обратных карт}
Следующий механизм родственен остаточным координатам
\cite{BD1993,EV2001,Lewis2024}; здесь записаны формулы, необходимые
для воспроизводимости.

\begin{proposition}[Достаточный критерий переноса]\label{prop:transfer-app2}
Пусть $A=(x,H,W)$ и $A^{-1}=(x,F,G)$ -- полиномиальные взаимно
обратные карты над $k$. Пусть $0\ne m\in k[x]$, $p\in k[x,y]$ и
\[
 \beta=(x,p(x,y)+m(x)z,y).
\]
Если $c(y)=\overline{H(x,p(x,y),y)}$ -- автоморфизм
$(k[x]/(m))[y]$, то $H\circ\beta$ дополняется до относительного
полиномиального автоморфизма.
\end{proposition}
\begin{proof}
Пусть $\psi$ -- обратный к $c$ в кольце $B[T]$, где $B=k[x]/(m)$.
Выберем подъём полинома
\[
 \varphi(x,T)\equiv W(x,p(x,\psi(T)),\psi(T))\pmod m.
\]
Тогда
\begin{equation}\label{eq:transfer-forward}
 A_{\rm new}=
 \left(x,H\circ\beta,
 \frac{W\circ\beta-\varphi(x,H\circ\beta)}m\right).
\end{equation}
Числитель последней компоненты обращается в нуль по модулю $m$.
Для входа $(x,T,V)$ положим
$S=mV+\varphi(x,T)$, $Y=F(x,T,S)$, $Z=G(x,T,S)$. Обратная карта:
\begin{equation}\label{eq:transfer-inverse}
 A_{\rm new}^{-1}(x,T,V)=\left(x,Z,\frac{Y-p(x,Z)}m\right).
\end{equation}
В $B[T]$ из обратимости $A$ получаем
$Y=p(x,\psi(T))$, $Z=\psi(T)$; деление снова точное.
После обращения $m$ формулы представляют композиции явно обратимых
карт. Инъективность локализации переносит обе тождественные композиции
в $k[x,y,z]$, включая все слои $m=0$.
\end{proof}

Условие на $c=c_0+c_1y+\sum_{j\ge2}c_jy^j$ эквивалентно обратимости
$c_1$ и нильпотентности всех $c_j$, $j\ge2$, в $B$.
Если $J=(c_j:j\ge2)$, $J^e=0$, обратный находится конечной итерацией
\[
 \psi_0=c_1^{-1}(T-c_0),\qquad
 \psi_{i+1}=c_1^{-1}\left(T-c_0-\sum_{j\ge2}c_j\psi_i^j\right).
\]
Разности последовательных итераций получают всё более высокие степени
$J$; после $e$ шагов достигается неподвижная точка. Таким образом,
вычисления должны вестись по полному $m$, а не только по его радикалу.

\subsection{Гармоническая модель и семейство $C_N$}\label{subsec:family-app2}
Тождество
\[
 D_3D_7-D_5^2+2=D_4
\]
показывает, что $\beta_1=(x,xz-y^2+2,y)$ переводит $\gamma_H$ в
$\gamma_0$. Поскольку $H(0,Y,Z)=-Z$, центр переноса равен $-y$
по модулю $x$, и предложение~\ref{prop:transfer-app2} применимо.
Можно взять поправку
\[
 \varphi_1(T)=42(a-1)T^6+175(1-a)T^4+(147a-161)T^2+8.
\]
Полагая $H_1=H\circ\beta_1$,
$W_1=(W\circ\beta_1-\varphi_1(H_1))/x$, имеем
\[
 W_1\gamma_H=b(42t^3-49t),\quad
 \Phi_H=(H_1,x-H_1^3+3H_1,W_1-b(42H_1^3-49H_1)).
\]
Степени компонент равны $(7,21,41)$. Обратное задаётся
\eqref{eq:transfer-inverse} и внешними обратными сдвигами.

Теперь пусть $N\ge1$. В $\Z[x,y]/(x^N)$ определим
\[
 \nu_0=-y,\qquad \nu_{i+1}=-y+x(\nu_i^2-2),\quad 0\le i<N,
 \qquad p_N\equiv2+x\nu_N-\nu_N^2\pmod{x^N},
\]
выбирая остаток $p_N$ степени по $x$ меньше $N$. Имеем
$p_N\equiv2-y^2\pmod x$. Положим
\begin{equation}\label{eq:family-CN}
 C_N(t)=\left(D_3,D_5,
 \frac{D_4-p_N(D_3,D_5)}{D_3^N}\right),\qquad
 \beta_N=(x,p_N+x^Nz,y).
\end{equation}
Деление в параметризации точное. Действительно,
$D_5=D_3(t^2-2)-t$ и $D_4=D_3t-t^2+2$.
Ошибка итерации $\nu_i(D_3,D_5)-t$ при каждом шаге получает дополнительный
множитель $D_3$. Следовательно, $\nu_N(D_3,D_5)\equiv t\pmod{D_3^N}$,
откуда и следует делимость.

Центр $H(x,p_N,y)$ равен $-y$ по модулю $x$. Поэтому по модулю $x^N$
его линейный коэффициент обратим, а все старшие нильпотентны.
Предложение~\ref{prop:transfer-app2} строит $A_N=(x,H_N,W_N)$ с
$H_N=H\circ\beta_N$ и $H_NC_N=t$. При $w_N=W_NC_N$ получаем
\begin{equation}\label{eq:family-rectifier}
 \Phi_N=(H_N,x-D_3(H_N),W_N-w_N(H_N)),\qquad
 \Phi_NC_N=(t,0,0).
\end{equation}
Тем самым $C_N$ действительно задаёт замкнутое вложение.
При $N=1$ можно выбрать указанный выше перенос и получить
$C_1=\gamma_H$, $\Phi_1=\Phi_H$. Различные представители поправки
могут менять компаньон ручным сдвигом; все последующие утверждения
относятся к построенным таким образом картам.

\subsection{Интегральность формул}
Все базовые коэффициенты лежат в
$\mathcal O=\Z[a]/(a^2-a+1)$. Конгруэнции для $P_x,F_x$ верны
уже над $\mathcal O$, а $b^{-1}=a$. Деление на $bx^5$ состоит из
умножения на $a$ и точного деления на моном $x^5$.
Рекурсии для $p_N$ и обратного центра не вводят числовых знаменателей:
линейный коэффициент центра равен $-1$ по модулю $x$ и обращается
конечным геометрическим рядом по модулю $x^N$.
Поэтому полиномиальные карты и их обратные специализируются в любое поле,
содержащее корень $a^2-a+1$, включая характеристики $2$ и $3$.
Это утверждение о формулах; доказательства дикости ниже относятся к
характеристике нуль. Степени после специализации могут уменьшаться.
Для $\mathbb F_q$ корень существует тогда и только тогда, когда
$\operatorname{char}\mathbb F_q=3$ или $q\equiv1\pmod3$:
в характеристике, отличной от $3$, элемент $-a$ является нетривиальным
кубическим корнем единицы. Отсутствие корня не исключает иных выпрямителей.
